\chapter{Fetch stage and branch prediction}
\label{ch:fetch}

\section{Fetch unit}

\file{hdl/fetch\_unit.v} owns the program counter and the \sig{ibus} AXI4-Lite
read master. It keeps at most one read in flight and raises the next \sig{ARVALID}
in the same cycle the current \sig{R} beat is accepted, so a latency-1 memory
sustains one instruction per cycle.

\subsection{Outstanding-transaction state}

The fetch unit has no state register. Its behaviour is encoded in four
independent flags:

\begin{table}[H]
\centering
\caption{Fetch unit state flags (\file{hdl/fetch\_unit.v})}
\label{tab:fetchflags}
\begin{tabular}{@{}lL{0.63\linewidth}@{}}
\toprule
\textbf{Flag} & \textbf{Meaning and update} \\
\midrule
\sig{ar\_pending\_q} & \sig{ARVALID} is up but has not been accepted. Set when an issue is
  not taken in its own cycle, cleared by \sig{ARREADY}. \\
\sig{inflight\_q}    & The address was accepted and the \sig{R} beat is still owed. Set on the
  AR handshake, cleared when the beat lands. Drives \sig{RREADY} directly. \\
\sig{discard\_q}     & The in-flight fetch is wrong-path. Set when a redirect arrives while the
  slot is busy, cleared when the beat is swallowed. \\
\sig{hold\_valid\_q} & The holding register contains an instruction that S2 could not accept.
  Set when a good beat lands while S2 is not ready, cleared on consumption or redirect. \\
\bottomrule
\end{tabular}
\end{table}

Although the flags are independent, the AXI slot they describe moves through a
well-defined cycle. Figure~\ref{fig:fetchfsm} draws that cycle as the four flags
produce it. It describes behaviour; no corresponding state register exists.

\begin{figure}[H]
\centering
\begin{tikzpicture}[x=1mm,y=1mm,>={Stealth[length=2.4mm]},semithick]
  % ---- the AXI read slot: three states -------------------------------
  \node[font=\scriptsize\itshape,anchor=west,color=accent] at (-14,16)
        {the AXI read slot};
  \node[stbi] (idle) at (0,0)    {IDLE};
  \node[stb]  (arp)  at (58,0)   {AR\_PEND};
  \node[stb]  (infl) at (29,-32) {IN\_FLIGHT};

  \draw[->,draw=accent] (idle) to[bend left=13]
      node[lbl,above,align=center,pos=0.5]
      {issue\_now \&\ !ARREADY} (arp);
  \draw[->,draw=accent,loop above,looseness=7,out=65,in=115] (arp) to
      node[lbl,above]{!ARREADY} (arp);
  \draw[->,draw=accent] (arp) to[bend left=13]
      node[lbl,right=1mm,align=left,pos=0.45]{ARREADY} (infl);
  \draw[->,draw=accent] (infl) to[bend left=13]
      node[lbl,left=1.5mm,align=right,pos=0.80]
      {RVALID\\\scriptsize beat consumed} (idle);
  \draw[->,draw=accent] (idle) to[bend left=13]
      node[lbl,right=1.5mm,align=left,pos=0.74]
      {issue\_now\\\scriptsize \&\ ARREADY} (infl);

  % ---- discard_q: a flag, not a fourth state -------------------------
  \node[font=\scriptsize\itshape,anchor=west,color=accent] at (-14,-50)
        {\sig{discard\_q}, orthogonal to the slot above};
  \node[stbi] (dc) at (6,-64)  {\sig{discard\_q} = 0};
  \node[stb]  (ds) at (62,-64) {\sig{discard\_q} = 1};

  \draw[->,draw=accent] (dc) to[bend left=17]
      node[lbl,above,align=center]{redirect\_i while the slot is busy} (ds);
  \draw[->,draw=accent] (ds) to[bend left=17]
      node[lbl,below,align=center]{RVALID --- the stale beat is swallowed} (dc);
\end{tikzpicture}
\caption{Instruction-fetch behaviour, derived from \sig{ar\_pending\_q},
\sig{inflight\_q} and \sig{discard\_q}. Double outlines mark the reset
condition. \sig{discard\_q} is drawn separately because it is exactly that in
the RTL: an independent flag, armed by a redirect from \emph{either} busy state,
because a read cannot be cancelled once \sig{ARVALID} is up. Neither machine
exists as a state register.}
\label{fig:fetchfsm}
\end{figure}

\subsection{Holding register}

When S2 stalls, a returned instruction parks in a one-entry holding register and
no new fetch is issued until it drains. S1 therefore can never run more than one
instruction ahead of S2, and the effective PC freezes for the duration of the
stall. The same mechanism is what silences the instruction bus during
\sig{WFI}: once the holding register is full, \sig{issue\_now} is low and no AR
goes out.

\subsection{Redirect with a read in flight}

An issued AXI read cannot be cancelled. When a redirect arrives while the slot is
busy, \sig{discard\_q} is set; the stale beat is accepted and dropped, and only
then does fetching restart at \sig{redirect\_pc}. The deferred address is held in
\sig{npc\_q} for the cycles in which nothing is being delivered.

\subsection{Fetch errors and reset behaviour}

A fetch that returns \sig{SLVERR} or \sig{DECERR} does not stall or trap in S1.
The word is delivered with \sig{fetch\_fault} set and traps as an instruction
access fault only if it reaches S2. A faulting fetch on a mispredicted path is
therefore flushed silently, which is the intended behaviour.

\sig{ARVALID} is explicitly gated with \sig{rst\_n\_i}:

\begin{center}
\ttfamily\small
assign ibus\_arvalid\_o = rst\_n\_i \& (ar\_pending\_q | issue\_now);
\end{center}

\begin{rtlnote}
The gate is not cosmetic. \sig{issue\_now} is combinational free-slot logic, so
without it the fetch unit requests \sig{RESET\_PC} while reset is still asserted.
AXI4-Lite forbids \sig{VALID} during reset, and a slave holding \sig{READY}
high in reset would accept a phantom address.
\end{rtlnote}

\section{Branch predictor}

\file{hdl/branch\_predictor.v} implements a combined BHT and BTB with an optional
return-address stack. Sizes come from the \sig{BP\_ENTRIES} and \sig{RAS\_DEPTH}
parameters forwarded by \file{cpu\_top}.

\subsection{Table organisation}

\begin{table}[H]
\centering
\caption{Predictor table entry, 128 entries by default}
\label{tab:btb}
\begin{tabular}{@{}llL{0.5\linewidth}@{}}
\toprule
\textbf{Field} & \textbf{Source} & \textbf{Purpose} \\
\midrule
\sig{valid}  & set on first update & Entry has been learned. \\
\sig{tag}    & \sig{PC[31:9]}      & Full remaining PC, so aliasing is detected rather than mispredicted. \\
\sig{state}  & 1 bit               & Last outcome of this branch; the 1-bit saturating counter the brief requires. \\
\sig{isret}  & 1 bit               & Entry was learned from a return; predict from the RAS top instead of the stored target. \\
\sig{target} & 32 bits             & Last resolved target. \\
\bottomrule
\end{tabular}
\end{table}

Indexing is direct-mapped on \sig{PC[8:2]}. Storing the target is not optional:
at fetch time the immediate has not been decoded, so a direction bit alone could
not redirect anything. A miss predicts not-taken, which is the only possibility
--- with no entry there is no stored target to jump to.

Updates happen only when a branch or jump commits in S2. Nothing in the predictor
is ever speculative, so there is no repair path that could itself be wrong after
a flush. Entries survive traps: trap and \sig{MRET} redirects are architectural,
come from \reg{mtvec}/\reg{mepc}, and bypass the predictor entirely.

\subsection{Return-address stack}

The BTB stores the last target, which is systematically wrong for a return called
from several sites. The RAS corrects that class of misprediction and is beyond
what the brief required.

\begin{table}[H]
\centering
\caption{RAS behaviour (\file{hdl/branch\_predictor.v})}
\label{tab:ras}
\begin{tabular}{@{}lL{0.62\linewidth}@{}}
\toprule
\textbf{Event} & \textbf{Action} \\
\midrule
Committed call  & Push PC+4. Identified by \sig{rd} $\in$ \{\sig{x1}, \sig{x5}\}, the register-hint
                  convention from the ISA specification. \\
Committed return & Pop. Identified by \sig{rs1} $\in$ \{\sig{x1}, \sig{x5}\} on \sig{JALR},
                  excluding the both-link case. \\
Both at once    & The popped slot is immediately refilled, so the top entry is replaced rather
                  than the depth changing. \\
Overflow        & The stack pointer wraps. \\
Underflow       & Prediction falls back to the stored BTB target. \\
\bottomrule
\end{tabular}
\end{table}

Overflow and underflow cost accuracy only, never correctness, because the
architectural answer always comes from \file{branch\_unit.v}. The pointer
arithmetic is sized to the pointer width rather than to the parameter, so the
wrap comparison stays exact for any legal \sig{RAS\_DEPTH}.

\subsection{Misprediction detection}

The prediction travels with the instruction through IF/DX and is compared in S2
against the real direction \emph{and} the real target. Either mismatch is a
misprediction. The hazard unit then asserts \sig{redirect}, the wrong-path fetch
in S1 is dropped, IF/DX takes a bubble, and the correct target is fetched the
next cycle. The branch instruction itself commits normally. Total cost: one
cycle.
